\section{Model}
\label{sec:model}
In this Section, I develop a model of fertility and housing decisions. Agents live a finite number of periods in an overlapping generations economy. They are subject to persistent idiosyncratic income shock and transitory fertility and housing taste shock. They derive utility from consumption, housing, and fertility, and can save using a one-period liquid asset or by purchasing a home. They can either rent or own a home, and pay transaction costs when seeling the home. Owners can borrow against their home equity, subject to liquidity constraint. Households' have a fertile period of their lives during which they can attempt to have kids, which give them utility but affect their housing and consumption needs. They also have a retirement age, after which they do not work and receive pension benefits. Such benefits are financed by a payroll tax which is imposed on workers, who supply labor inelastically. 

\subsection{The economy}
I begin by describing agents and their preferences, the income process, the markets in which agents can trade and the fertility process. Time is discrete.\\

\subsubsection*{Households and Preferences.}
The economy is populated by overlapping generation of households, who live for at most T periods. Households are also denoted as `adults'. Adult age is denoted by $a\in[0,J]$. Agents can die at each period, and I denote $s_a$ their probability of surviving to next period.  Households work before retirement age $a_R$ and can attempt births during the
fertile window $a\in[a_F^0,a_F^1]$. Let $n_t$ denote children ever born and $m_t$ children currently
at home. Children are dependents and do not work. Each dependent child matures
and leaves the household with a fixed probability \textcolor{red}{fix} each period, independently
of the other children. Maturation reduces $m_t$ but does not change $n_t$. Households value consumption, housing services, and children at home, an discount future utility by $\beta$. At any period $t$, agents have flow utility

\begin{equation*}
	u_t(c_t,s_t;m_t)
	=\frac{\mathcal C(c_t,s_t;m_t)^{1-\sigma}}{1-\sigma}
	+\psi_t m_t.
\end{equation*}
The parameter $\sigma>0$ governs curvature over the consumption aggregator
$\mathcal C$, with the logarithmic limit understood when $\sigma=1$, while
$\psi_t$ measures the direct utility benefit of a child at home.
For a given number of children, $\mathcal C$ is increasing and concave in
consumption and housing services. Children increase the household's needs,
so $\mathcal C_m<0$, and raise the marginal contribution of housing services
to the aggregator, $\mathcal C_{sm}>0$. The functional form of
$\mathcal C$ is specified in the quantification.

Households either rent or own their residence. An owner occupies $h_t$ units
of housing, while a renter occupies $h_t^R$ units. Housing services are
\begin{equation}
 s_t=\mathbf{1}\{\text{own}\}\chi h_t
     +\mathbf{1}\{\text{rent}\}h_t^R,
 \qquad \chi>1.
 \label{eq:env-housing-services}
\end{equation}
The parameter $\chi$ captures the additional services households obtain from
owning a home of a given size.

At death, households value the estate they leave. Let $w^e_{t+1}$ denote
bequeathable wealth, comprising liquid wealth and housing value. Bequest
utility is
\begin{equation}
 B(w^e_{t+1})
 =\theta_0\frac{(\theta_1+w^e_{t+1})^{1-\sigma}}{1-\sigma}.
 \label{eq:env-bequest-utility}
\end{equation}
The parameter $\theta_0$ scales the strength of the bequest motive, while
$\theta_1$ shifts the wealth level at which it is evaluated; the logarithmic
limit again applies when $\sigma=1$.

\paragraph{Fertility.}
During fertile ages, a household chooses whether to attempt a birth. An
attempt succeeds with age-specific probability $\pi_a$ and, upon success,
increases both $n_t$ and $m_t$ by one. The first birth carries a one-time
utility cost $\xi$. Households observe extreme-value taste shocks
$\varepsilon_t^F$ before choosing whether to attempt; their scale is
$\kappa_1$ for the first birth and $\kappa_C$ for subsequent births.
Birth success is then realized. Households subsequently observe housing
taste shocks $\varepsilon_t^H$ and choose housing and saving conditional on
the realized number of children. Thus, housing can adjust to a birth within
the period.

\paragraph{Production and income.}
Competitive firms produce the consumption good using effective labor $L_t$
with constant productivity $A$:
\begin{equation}
 Y_t=A L_t,\qquad w=A.
 \label{eq:env-production}
\end{equation}
The wage per unit of effective labor therefore equals labor productivity.
A working household supplies $e_a z_t$ units of effective labor, where $e_a$
is the age--earnings profile and $z_t$ exposes the household to idiosyncratic
earnings risk. Earnings are taxed at rate $\tau^w$. Retirees receive a
pay-as-you-go pension $\varpi_t$, financed by current payroll taxes.
Disposable income is
\begin{equation}
 y_t^d(a,z_t)=
 \begin{cases}
 (1-\tau^w)w e_a z_t, & a<a_R,\\
 \varpi_t, & a\geq a_R.
 \end{cases}
 \label{eq:env-income}
\end{equation}
Households also receive a lump-sum transfer $T_t$ from property-tax revenue,
which is rebated equally across households.

\paragraph{Housing markets.}
A unit of housing trades at price $P_t$. Aggregate housing supply is
\begin{equation}
 H_t^s=H_0 P_t^{\eta}.
 \label{eq:env-housing-supply}
\end{equation}
Here $H_0$ sets the scale of supply and $\eta$ is its price elasticity.
Owned housing can be as large as $\bar h$, whereas rental housing is limited
to $\bar h^R<\bar h$. The largest homes are therefore available only through
ownership. Housing depreciates at rate $\delta_H$ and is subject to a
property tax at rate $\tau_t^p$. Selling an owner-occupied home costs a
fraction $\psi^s$ of its sale value.

Competitive landlords rent housing at a unit rent $r_t$. A landlord purchases
at $P_t$, pays depreciation and property tax, and resells at $P_{t+1}$.
Equating the return to the gross bond return $R_b$ gives
\begin{equation}
 r_t=(R_b+\delta_H+\tau_t^p)P_t-P_{t+1}.
 \label{eq:env-rent}
\end{equation}
Rent covers financing, depreciation, and property taxes, net of the change
in the value of the house.

\paragraph{Saving and mortgage finance.}
Households save in a one-period bond with gross return $R_b$. Borrowing is
secured by housing: an owner can borrow up to a fraction $\phi$ of the
home's purchase value, while renters cannot borrow. Buying a home requires
a down payment of $(1-\phi)P_t h_t$. The household must finance this payment
out of beginning-of-period liquid wealth and the proceeds from selling its
previous home, net of selling costs. Current earnings and transfers arrive
after the housing transaction and cannot finance the down payment.
Households that retain their home do not make a new down payment.
